\documentclass[UTF8,12pt,letterpaper]{ctexart}

\usepackage{amsmath}
\usepackage{graphicx}
\graphicspath{{images/}}

\title{我的第一份 \LaTeX{} 文档}
\author{
  Nicolas%
  \thanks{Funded by the College of Electrical and Automation Engineering,
  Nanjing Normal University}
}
\date{2026年7月14日}

\begin{document}

\maketitle

\section{文本格式}

We have now added a title, author, and date to our first
\LaTeX{} document!

% 这一行是注释，不会显示在最终文档中。

Some of the \textbf{greatest}
\emph{discoveries} in \underline{science}
were made by \textbf{\textit{accident}}.

\section{数学公式}

Einstein's mass--energy equivalence is

\[
E = mc^2.
\]

另一个公式为

\[
d = \frac{k\varphi(n)+1}{e}.
\]

\section{表格}

\begin{table}[htbp]
  \centering
  \caption{一个简单的数字表格}
  \label{tab:number-example}

  \begin{tabular}{|c|c|c|}
    \hline
    1 & 2 & 3 \\
    \hline
    4 & 5 & 6 \\
    \hline
    7 & 8 & 9 \\
    \hline
  \end{tabular}
\end{table}

表~\ref{tab:number-example} 展示了一个简单的 $3\times3$ 数字表格。

\end{document}
